\itemhead{复习题二 · 第17题（教材p.63）}{XJ-B1-C02-RV-P063-Y1035}
\label{XJ-B1-C02-RV-P063-Y1035}
{\footnotesize 湘教版教材p.63 / PDF第70页 / 复习题二 / 第17题}
\field{原题}{为了在夏季降温和冬季供暖时减少能源损耗，房屋的屋顶和外墙需要建造隔热层。某幢建筑物要建造可使用 20 年的隔热层，每厘米厚的隔热层建造成本为 6 万元。该建筑物每年的能源消耗费用 $C$（万元）与隔热层厚度 $x$（$\mathrm{cm}$）满足如下关系式：$C=\dfrac k{3x+5}\ (0\le x\le10)$。若不建隔热层，每年能源消耗费用为 8 万元。设 $y$（万元）为隔热层建造费用与 20 年的能源消耗费用之和。
\subq{1}求 $k$ 的值及 $y$ 的表达式。
\subq{2}隔热层修建多厚时，总费用 $y$ 最小？并求最小值。}
\field{参考答案}{（1）$k=40,y=6x+800/(3x+5)$，$0\le x\le10$；（2）置 $t=3x+5$，$y=2t+800/t-10\ge70$，$t=20$ 即 $x=5$ cm 时最小70万元。}
{\small \field{关联知识}{I03 基本不等式、等号与最值建模（教材帮p.52、53、54 / PDF61、62、63）；F03 函数求值、解析式与函数建模（教材帮p.80、87、88、89、90、131 / PDF89、96、97、98、99、140）}
\field{方法与易错点}{列正数条件与固定量，构造基本不等式，解等号条件并代回实际限制。 易错：把一个下界误当最小值；多个等号条件不能同时实现；漏单位和整数限制。}
\field{答案依据}{答案册相关题组：PDF第15、16页（册内页14、15）；按当前题库原题逐项重算，不按旧版题号机械对套。}}
